% --- Archon physics formalization source begin ---
% archon:physics
% archon:covers PhyXMiniProblems/problem_phyx_mini_0603.lean
% archon:source-report reports/phyx_mini/problem_phyx_mini_0603.source.json

\chapter{Physics problem phyx\_mini\_0603}
\label{ch:PhyXMiniProblems_problem_phyx_mini_0603}

\paragraph{Problem source.}
\#\# Physical scenario

In a monovalent metal, one electron per atom is free to roam throughout the object. Evidently the energy of the composite structure must be less than the energy of the isolated atoms. This problem offers a crude but illuminating explanation for the cohesiveness of metals.

\#\# Question

When these atoms come together to form a metal, we get \$N\$ electrons in a much larger infinite square well of width \$Na\$. Because of the Pauli exclusion principle (which we will discuss in Chapter 5) there can only be one electron (two, if you include spin, but let's ignore that) in each allowed state. What is the lowest energy for this system?

\#\# Answer choices

- A: A: \textbackslash{}frac\{N \textbackslash{}pi\textasciicircum{}2 \textbackslash{}hbar\textasciicircum{}2\}\{2m\_e (Na)\textasciicircum{}2\}

- B: B: \textbackslash{}frac\{\textbackslash{}pi\textasciicircum{}2 \textbackslash{}hbar\textasciicircum{}2\}\{2m\_e a\textasciicircum{}2\} \textbackslash{}cdot \textbackslash{}frac\{N\textasciicircum{}2\}\{6\}

- C: C: \textbackslash{}frac\{\textbackslash{}pi\textasciicircum{}2 \textbackslash{}hbar\textasciicircum{}2\}\{2m\_e a\textasciicircum{}2\} \textbackslash{}cdot \textbackslash{}frac\{(N + 1)(2N + 1)\}\{6N\}

- D: D: \textbackslash{}frac\{N \textbackslash{}pi\textasciicircum{}2 \textbackslash{}hbar\textasciicircum{}2\}\{2m\_e a\textasciicircum{}2\}

\#\# Figure caption (auxiliary; use the image as primary evidence)

The image appears to be a diagram featuring a straight horizontal line with several equally spaced open circles positioned along it. These circles likely represent elements or points of interest on the line. The line is bounded by two perpendicular lines, forming a rectangular area.

Below the line, there is a label and an arrow indicating the horizontal length. The text “\textbackslash{}(Na\textbackslash{})” (with \textbackslash{}(N\textbackslash{}) in italics and \textbackslash{}(a\textbackslash{}) in a different font style) is placed beneath the arrow, suggesting it represents a distance or measurement.

The diagram includes a series of circles, with one section having more closely spaced circles and another section with circles spread out. The text "..." is placed between two circles, implying a continuation or an unspecified number of additional points.

Overall, this structure likely represents a simplified schematic or conceptual illustration, possibly from a scientific or mathematical context.

\#\# Dataset metadata

Quantum Phenomena; Physical Model Grounding Reasoning, Numerical Reasoning

\paragraph{Recorded answer/context.}
C: C: \textbackslash{}frac\{\textbackslash{}pi\textasciicircum{}2 \textbackslash{}hbar\textasciicircum{}2\}\{2m\_e a\textasciicircum{}2\} \textbackslash{}cdot \textbackslash{}frac\{(N + 1)(2N + 1)\}\{6N\}

\paragraph{Figure/image path.}
/root/proposal\_for\_physic/science-mango/phyx\_mini\_run/phyx\_data/test\_image/603.png

\paragraph{Formalization target.}
create a compiling Lean file with sorry bodies at `PhyXMiniProblems/problem\_phyx\_mini\_0603.lean`. The Lean declarations must preserve the physical quantities, dimensions or dimensional roles, figure labels, governing-law hypotheses, and final relation expressed by this problem.
Use Mathlib/Physlib names found through LeanExplore where available. If a physics API is missing, introduce faithful local abstractions rather than scalar placeholder aliases.

\begin{theorem}[Physics formalization target]
\label{thm:physics:phyx_mini_0603:target}
\lean{PhyXMiniProblems.ProblemPhyXMini0603.problem_phyx_mini_0603}
\uses{def:physics:phyx-mini-0603:phyxminiproblems-problemphyxmini0603-energyinjoules, def:physics:phyx-mini-0603:phyxminiproblems-problemphyxmini0603-metalwellsetup, def:physics:phyx-mini-0603:phyxminiproblems-problemphyxmini0603-matchesmonovalentmetalscenario, def:physics:phyx-mini-0603:phyxminiproblems-problemphyxmini0603-matchessuppliedmetalwellfigure, def:physics:phyx-mini-0603:phyxminiproblems-problemphyxmini0603-hasphysicalmetalwellparameters, def:physics:phyx-mini-0603:phyxminiproblems-problemphyxmini0603-satisfiesinfinitesquarewellspectrum, def:physics:phyx-mini-0603:phyxminiproblems-problemphyxmini0603-satisfiesspinlessfermiongroundstate, def:physics:phyx-mini-0603:phyxminiproblems-problemphyxmini0603-singleatomenergyscaleinjoules, def:physics:phyx-mini-0603:phyxminiproblems-problemphyxmini0603-matchesdisplayedenergychoice}
The assigned autoformalize agent should translate this physics problem into Lean declarations in the covered file, with theorem and lemma proof bodies written as `by sorry`.
\end{theorem}
\begin{proof}
This is an autoformalization task, not a proof task. Produce statements that can later be proved by the physics prover without weakening the physical model.
\end{proof}

% --- Archon declaration topology begin ---
\section{Declaration topology}
\begin{definition}[Length Quantity]
  \label{def:physics:phyx-mini-0603:phyxminiproblems-problemphyxmini0603-lengthquantity}
  \lean{PhyXMiniProblems.ProblemPhyXMini0603.LengthQuantity}
  A nonnegative, unit-independent physical length
\end{definition}
\begin{proof}
  This is the declared model definition of Length Quantity.
\end{proof}
\begin{definition}[Mass Quantity]
  \label{def:physics:phyx-mini-0603:phyxminiproblems-problemphyxmini0603-massquantity}
  \lean{PhyXMiniProblems.ProblemPhyXMini0603.MassQuantity}
  A nonnegative, unit-independent physical mass
\end{definition}
\begin{proof}
  This is the declared model definition of Mass Quantity.
\end{proof}
\begin{definition}[action Dimension]
  \label{def:physics:phyx-mini-0603:phyxminiproblems-problemphyxmini0603-actiondimension}
  \lean{PhyXMiniProblems.ProblemPhyXMini0603.actionDimension}
  The dimension of action, mass * length\textasciicircum{}2 / time
\end{definition}
\begin{proof}
  This is the declared model definition of action Dimension.
\end{proof}
\begin{definition}[Action Quantity]
  \label{def:physics:phyx-mini-0603:phyxminiproblems-problemphyxmini0603-actionquantity}
  \lean{PhyXMiniProblems.ProblemPhyXMini0603.ActionQuantity}
  \uses{def:physics:phyx-mini-0603:phyxminiproblems-problemphyxmini0603-actiondimension}
  A nonnegative, unit-independent physical action such as reduced Planck action
\end{definition}
\begin{proof}
  This is the declared model definition of Action Quantity.
\end{proof}
\begin{definition}[length In Meters]
  \label{def:physics:phyx-mini-0603:phyxminiproblems-problemphyxmini0603-lengthinmeters}
  \lean{PhyXMiniProblems.ProblemPhyXMini0603.lengthInMeters}
  \uses{def:physics:phyx-mini-0603:phyxminiproblems-problemphyxmini0603-lengthquantity}
  Read a physical length in coherent SI units, metres
\end{definition}
\begin{proof}
  This is the declared model definition of length In Meters.
\end{proof}
\begin{definition}[mass In Kilograms]
  \label{def:physics:phyx-mini-0603:phyxminiproblems-problemphyxmini0603-massinkilograms}
  \lean{PhyXMiniProblems.ProblemPhyXMini0603.massInKilograms}
  \uses{def:physics:phyx-mini-0603:phyxminiproblems-problemphyxmini0603-massquantity}
  Read a physical mass in coherent SI units, kilograms
\end{definition}
\begin{proof}
  This is the declared model definition of mass In Kilograms.
\end{proof}
\begin{definition}[action In Joule Seconds]
  \label{def:physics:phyx-mini-0603:phyxminiproblems-problemphyxmini0603-actioninjouleseconds}
  \lean{PhyXMiniProblems.ProblemPhyXMini0603.actionInJouleSeconds}
  \uses{def:physics:phyx-mini-0603:phyxminiproblems-problemphyxmini0603-actionquantity}
  Read a physical action in coherent SI units, joule-seconds
\end{definition}
\begin{proof}
  This is the declared model definition of action In Joule Seconds.
\end{proof}
\begin{definition}[energy In Joules]
  \label{def:physics:phyx-mini-0603:phyxminiproblems-problemphyxmini0603-energyinjoules}
  \lean{PhyXMiniProblems.ProblemPhyXMini0603.energyInJoules}
  Read a physical energy in coherent SI units, joules
\end{definition}
\begin{proof}
  This is the declared model definition of energy In Joules.
\end{proof}
\begin{definition}[Particle Species]
  \label{def:physics:phyx-mini-0603:phyxminiproblems-problemphyxmini0603-particlespecies}
  \lean{PhyXMiniProblems.ProblemPhyXMini0603.ParticleSpecies}
  Particle species singled out by the problem statement
\end{definition}
\begin{proof}
  This is the declared model definition of Particle Species.
\end{proof}
\begin{definition}[Metal Valence Model]
  \label{def:physics:phyx-mini-0603:phyxminiproblems-problemphyxmini0603-metalvalencemodel}
  \lean{PhyXMiniProblems.ProblemPhyXMini0603.MetalValenceModel}
  Valence model for the metal atoms
\end{definition}
\begin{proof}
  This is the declared model definition of Metal Valence Model.
\end{proof}
\begin{definition}[Spin Treatment]
  \label{def:physics:phyx-mini-0603:phyxminiproblems-problemphyxmini0603-spintreatment}
  \lean{PhyXMiniProblems.ProblemPhyXMini0603.SpinTreatment}
  Whether electron spin degeneracy is retained in the occupancy count
\end{definition}
\begin{proof}
  This is the declared model definition of Spin Treatment.
\end{proof}
\begin{definition}[Confinement Model]
  \label{def:physics:phyx-mini-0603:phyxminiproblems-problemphyxmini0603-confinementmodel}
  \lean{PhyXMiniProblems.ProblemPhyXMini0603.ConfinementModel}
  Confining potential used for the delocalized electrons
\end{definition}
\begin{proof}
  This is the declared model definition of Confinement Model.
\end{proof}
\begin{definition}[Figure Feature]
  \label{def:physics:phyx-mini-0603:phyxminiproblems-problemphyxmini0603-figurefeature}
  \lean{PhyXMiniProblems.ProblemPhyXMini0603.FigureFeature}
  Literal graphical features visible in image 603.png
\end{definition}
\begin{proof}
  This is the declared model definition of Figure Feature.
\end{proof}
\begin{definition}[Figure Width Label]
  \label{def:physics:phyx-mini-0603:phyxminiproblems-problemphyxmini0603-figurewidthlabel}
  \lean{PhyXMiniProblems.ProblemPhyXMini0603.FigureWidthLabel}
  Semantic reading of the width annotation printed below the well
\end{definition}
\begin{proof}
  This is the declared model definition of Figure Width Label.
\end{proof}
\begin{definition}[Metal Well Figure]
  \label{def:physics:phyx-mini-0603:phyxminiproblems-problemphyxmini0603-metalwellfigure}
  \lean{PhyXMiniProblems.ProblemPhyXMini0603.MetalWellFigure}
  \uses{def:physics:phyx-mini-0603:phyxminiproblems-problemphyxmini0603-figurefeature, def:physics:phyx-mini-0603:phyxminiproblems-problemphyxmini0603-figurewidthlabel}
  Literal visual data transcribed from the supplied metal-well diagram
\end{definition}
\begin{proof}
  This is the declared model definition of Metal Well Figure.
\end{proof}
\begin{definition}[Metal Well Setup]
  \label{def:physics:phyx-mini-0603:phyxminiproblems-problemphyxmini0603-metalwellsetup}
  \lean{PhyXMiniProblems.ProblemPhyXMini0603.MetalWellSetup}
  \uses{def:physics:phyx-mini-0603:phyxminiproblems-problemphyxmini0603-lengthquantity, def:physics:phyx-mini-0603:phyxminiproblems-problemphyxmini0603-massquantity, def:physics:phyx-mini-0603:phyxminiproblems-problemphyxmini0603-actionquantity, def:physics:phyx-mini-0603:phyxminiproblems-problemphyxmini0603-particlespecies, def:physics:phyx-mini-0603:phyxminiproblems-problemphyxmini0603-metalvalencemodel, def:physics:phyx-mini-0603:phyxminiproblems-problemphyxmini0603-spintreatment, def:physics:phyx-mini-0603:phyxminiproblems-problemphyxmini0603-confinementmodel, def:physics:phyx-mini-0603:phyxminiproblems-problemphyxmini0603-metalwellfigure}
  Independent physical quantities in the model. In particular, groundStateEnergy is an unconstrained dimensionful field here: it is not defined from any answer choice or from the requested closed form
\end{definition}
\begin{proof}
  This is the declared model definition of Metal Well Setup.
\end{proof}
\begin{definition}[Matches Monovalent Metal Scenario]
  \label{def:physics:phyx-mini-0603:phyxminiproblems-problemphyxmini0603-matchesmonovalentmetalscenario}
  \lean{PhyXMiniProblems.ProblemPhyXMini0603.MatchesMonovalentMetalScenario}
  \uses{def:physics:phyx-mini-0603:phyxminiproblems-problemphyxmini0603-metalwellsetup}
  The prose model: a monovalent metal contributes one conduction electron per atom, spin is ignored, and the electrons occupy a one-dimensional infinite square well
\end{definition}
\begin{proof}
  This is the declared model definition of Matches Monovalent Metal Scenario.
\end{proof}
\begin{definition}[Matches Supplied Metal Well Figure]
  \label{def:physics:phyx-mini-0603:phyxminiproblems-problemphyxmini0603-matchessuppliedmetalwellfigure}
  \lean{PhyXMiniProblems.ProblemPhyXMini0603.MatchesSuppliedMetalWellFigure}
  \uses{def:physics:phyx-mini-0603:phyxminiproblems-problemphyxmini0603-lengthinmeters, def:physics:phyx-mini-0603:phyxminiproblems-problemphyxmini0603-metalwellsetup}
  Primary-image evidence: two walls and a floor bound the well, eight open circles are explicitly drawn with an ellipsis for continuation, and the double-headed span arrow is labelled N a. The last field is the calibrated physical reading of that label, not an energy conclusion
\end{definition}
\begin{proof}
  This is the declared model definition of Matches Supplied Metal Well Figure.
\end{proof}
\begin{definition}[Has Physical Metal Well Parameters]
  \label{def:physics:phyx-mini-0603:phyxminiproblems-problemphyxmini0603-hasphysicalmetalwellparameters}
  \lean{PhyXMiniProblems.ProblemPhyXMini0603.HasPhysicalMetalWellParameters}
  \uses{def:physics:phyx-mini-0603:phyxminiproblems-problemphyxmini0603-lengthinmeters, def:physics:phyx-mini-0603:phyxminiproblems-problemphyxmini0603-massinkilograms, def:physics:phyx-mini-0603:phyxminiproblems-problemphyxmini0603-actioninjouleseconds, def:physics:phyx-mini-0603:phyxminiproblems-problemphyxmini0603-metalwellsetup}
  Positivity conditions selecting a nondegenerate physical metal sample
\end{definition}
\begin{proof}
  This is the declared model definition of Has Physical Metal Well Parameters.
\end{proof}
\begin{definition}[Satisfies Infinite Square Well Spectrum]
  \label{def:physics:phyx-mini-0603:phyxminiproblems-problemphyxmini0603-satisfiesinfinitesquarewellspectrum}
  \lean{PhyXMiniProblems.ProblemPhyXMini0603.SatisfiesInfiniteSquareWellSpectrum}
  \uses{def:physics:phyx-mini-0603:phyxminiproblems-problemphyxmini0603-lengthinmeters, def:physics:phyx-mini-0603:phyxminiproblems-problemphyxmini0603-massinkilograms, def:physics:phyx-mini-0603:phyxminiproblems-problemphyxmini0603-actioninjouleseconds, def:physics:phyx-mini-0603:phyxminiproblems-problemphyxmini0603-energyinjoules, def:physics:phyx-mini-0603:phyxminiproblems-problemphyxmini0603-metalwellsetup}
  The one-dimensional infinite-square-well spectrum E\_n = n\textasciicircum{}2 * pi\textasciicircum{}2 * hbar\textasciicircum{}2 / (2 m L\textasciicircum{}2) for every positive quantum number. Physlib currently has no dedicated infinite-square-well spectrum declaration, so this is a local governing-law interface rather than the requested total-energy formula
\end{definition}
\begin{proof}
  This is the declared model definition of Satisfies Infinite Square Well Spectrum.
\end{proof}
\begin{definition}[Is Admissible Spinless Occupation]
  \label{def:physics:phyx-mini-0603:phyxminiproblems-problemphyxmini0603-isadmissiblespinlessoccupation}
  \lean{PhyXMiniProblems.ProblemPhyXMini0603.IsAdmissibleSpinlessOccupation}
  \uses{def:physics:phyx-mini-0603:phyxminiproblems-problemphyxmini0603-metalwellsetup}
  An admissible spinless occupation of N electrons is a finite set of exactly N positive quantum numbers. Use of a Finset expresses the Pauli rule that no level is occupied more than once when spin degeneracy is ignored
\end{definition}
\begin{proof}
  This is the declared model definition of Is Admissible Spinless Occupation.
\end{proof}
\begin{definition}[Satisfies Spinless Fermion Ground State]
  \label{def:physics:phyx-mini-0603:phyxminiproblems-problemphyxmini0603-satisfiesspinlessfermiongroundstate}
  \lean{PhyXMiniProblems.ProblemPhyXMini0603.SatisfiesSpinlessFermionGroundState}
  \uses{def:physics:phyx-mini-0603:phyxminiproblems-problemphyxmini0603-energyinjoules, def:physics:phyx-mini-0603:phyxminiproblems-problemphyxmini0603-metalwellsetup, def:physics:phyx-mini-0603:phyxminiproblems-problemphyxmini0603-isadmissiblespinlessoccupation}
  The ground-state governing principle. The setup's occupied levels form an admissible Pauli occupation, the total energy is additive over them, and that sum is minimal among every admissible occupation. This predicate contains no closed form for the minimum and no answer-choice value
\end{definition}
\begin{proof}
  This is the declared model definition of Satisfies Spinless Fermion Ground State.
\end{proof}
\begin{definition}[Answer Choice]
  \label{def:physics:phyx-mini-0603:phyxminiproblems-problemphyxmini0603-answerchoice}
  \lean{PhyXMiniProblems.ProblemPhyXMini0603.AnswerChoice}
  Labels attached to the four symbolic energy expressions in the source
\end{definition}
\begin{proof}
  This is the declared model definition of Answer Choice.
\end{proof}
\begin{definition}[single Atom Energy Scale In Joules]
  \label{def:physics:phyx-mini-0603:phyxminiproblems-problemphyxmini0603-singleatomenergyscaleinjoules}
  \lean{PhyXMiniProblems.ProblemPhyXMini0603.singleAtomEnergyScaleInJoules}
  \uses{def:physics:phyx-mini-0603:phyxminiproblems-problemphyxmini0603-lengthinmeters, def:physics:phyx-mini-0603:phyxminiproblems-problemphyxmini0603-massinkilograms, def:physics:phyx-mini-0603:phyxminiproblems-problemphyxmini0603-actioninjouleseconds, def:physics:phyx-mini-0603:phyxminiproblems-problemphyxmini0603-metalwellsetup}
  The single-atom ground-level energy scale pi\textasciicircum{}2 hbar\textasciicircum{}2 / (2 m\_e a\textasciicircum{}2), expressed in joules
\end{definition}
\begin{proof}
  This is the declared model definition of single Atom Energy Scale In Joules.
\end{proof}
\begin{definition}[displayed Energy In Joules]
  \label{def:physics:phyx-mini-0603:phyxminiproblems-problemphyxmini0603-displayedenergyinjoules}
  \lean{PhyXMiniProblems.ProblemPhyXMini0603.displayedEnergyInJoules}
  \uses{def:physics:phyx-mini-0603:phyxminiproblems-problemphyxmini0603-lengthinmeters, def:physics:phyx-mini-0603:phyxminiproblems-problemphyxmini0603-massinkilograms, def:physics:phyx-mini-0603:phyxminiproblems-problemphyxmini0603-actioninjouleseconds, def:physics:phyx-mini-0603:phyxminiproblems-problemphyxmini0603-metalwellsetup, def:physics:phyx-mini-0603:phyxminiproblems-problemphyxmini0603-answerchoice, def:physics:phyx-mini-0603:phyxminiproblems-problemphyxmini0603-singleatomenergyscaleinjoules}
  The four symbolic energy expressions printed beside choices A--D
\end{definition}
\begin{proof}
  This is the declared model definition of displayed Energy In Joules.
\end{proof}
\begin{definition}[recorded Dataset Answer]
  \label{def:physics:phyx-mini-0603:phyxminiproblems-problemphyxmini0603-recordeddatasetanswer}
  \lean{PhyXMiniProblems.ProblemPhyXMini0603.recordedDatasetAnswer}
  \uses{def:physics:phyx-mini-0603:phyxminiproblems-problemphyxmini0603-answerchoice}
  Dataset metadata recording answer C; it is not used as a theorem premise
\end{definition}
\begin{proof}
  This is the declared model definition of recorded Dataset Answer.
\end{proof}
\begin{definition}[Matches Displayed Energy Choice]
  \label{def:physics:phyx-mini-0603:phyxminiproblems-problemphyxmini0603-matchesdisplayedenergychoice}
  \lean{PhyXMiniProblems.ProblemPhyXMini0603.MatchesDisplayedEnergyChoice}
  \uses{def:physics:phyx-mini-0603:phyxminiproblems-problemphyxmini0603-energyinjoules, def:physics:phyx-mini-0603:phyxminiproblems-problemphyxmini0603-metalwellsetup, def:physics:phyx-mini-0603:phyxminiproblems-problemphyxmini0603-answerchoice, def:physics:phyx-mini-0603:phyxminiproblems-problemphyxmini0603-displayedenergyinjoules}
  The physical ground-state energy agrees with one displayed expression
\end{definition}
\begin{proof}
  This is the declared model definition of Matches Displayed Energy Choice.
\end{proof}
% --- Archon declaration topology end ---
% --- Archon physics formalization source end ---
